平面ベクトル$\va*{p}$，$\va*{q}$の内積を$\va*{p}\cdot\va*{q}$と表す．
$f$は平面上の一次変換とする．
\begin{enumerate}
  \item $\va*{p}$，$\va*{q}$がたがいに直交する単にベクトルとすると，$T=f(\va*{p})\cdot\va*{p}+f(\va*{q})\cdot\va*{q}$は，ベクトルの組$\va*{p}$，$\va*{q}$の取り方によらないで，$f$によって決まる値であることを示せ．
  \item 原点$\mathrm{O}$を通る$2$つの定直線$l$と$m$があって，$f$によって$l$上の任意の点$\mathrm{R}$は$\mathrm{R}$自身に移され，$m$上の任意の点$\mathrm{S}$は$\mathrm{OS}$の中点$\mathrm{S}'$に移されるとする．このとき$f$に対する$T$の値を求めよ．
\end{enumerate}
